\section{Introduction}

A predictive benchmark is a measurement design, not a neutral ruler. In chemometrics, the meaning of a held-out score depends on how the validation sample represents the population to which a model will be applied; resampling choices can alter extrapolation difficulty, uncertainty, and the scientific claim supported by the score \citep{esbensen2010proper,tropsha2003earnest,gramatica2007principles,lasfar2024robustness,kiraly2025leakage,camacho2026validation}. Molecular property prediction makes this dependence especially visible because the same molecular universe can be partitioned by observation, structural framework, chronology, or explicit leakage constraints.

MoleculeNet popularized common molecular-property tasks and scaffold-based evaluation \citep{wu2018moleculenet}. Bemis--Murcko splitting \citep{bemis1996properties} is widely used to reduce direct framework overlap, yet scaffold-disjoint partitions can still differ in test size, response prevalence or mean, scaffold concentration, and the particular chemical series withheld. Comparative studies have documented substantial variation across split schemes and random seeds \citep{yang2019learned,deng2023systematic,li2026partition}; leakage-aware and prospective-proxy approaches such as DataSAIL and SIMPD address related but distinct deployment questions \citep{joeres2025datasail,landrum2023simpd}. Dataset identity and preprocessing rules are likewise part of the QSAR claim rather than interchangeable implementation details \citep{nael2026dataset}.

This study asks a narrower controlled question: \emph{what happens when a scaffold-disjoint partition is selected to have a lower train--test target-mean mismatch, while test-set size and candidate-search budget are held fixed?} This wording is deliberate. The intervention changes the selected molecular subset and therefore can also change structural composition; it does not ``isolate'' the full target distribution as a causal variable. Instead, it quantifies the performance consequence of a target-mean-aware selection rule inside a fixed target-blind candidate space.

For each dataset and each of 20 pre-specified partition seeds, we generate a fixed-budget pool of target-blind scaffold candidates. A size-matched baseline is selected using test-size deviation only. A target-mean-balanced counterpart is then chosen from candidates with \emph{exactly the same test-set size} by minimizing the absolute train--test target-mean gap. Candidate generation, test cardinality, scaffold rule, and search budget are therefore controlled at the protocol level, while collateral changes in the selected test-set composition are measured rather than assumed away.

Search effort is treated as part of the benchmark definition. A larger candidate pool provides more opportunities to find an extreme low-gap partition, creating a researcher degree of freedom analogous to expanding a hyperparameter search. Candidate budgets are therefore audited and frozen before model outcomes are examined. Every materialized partition is stored in a molecule-level manifest and identified by a cryptographic hash, and the inferential unit is the unique paired partition rather than a model seed.

Two chemically meaningful protocol sensitivities test whether the observed conclusion survives alternative meanings of molecular structure. First, an acyclic molecule has an empty Bemis--Murcko framework. We compare a \emph{single-group} convention, in which all acyclic molecules share one scaffold identity, with a \emph{singleton} convention, in which each acyclic molecule receives its own identity. Second, disconnected source records can contain salts, counterions, complexes, or mixtures. A deterministic dominant-fragment mapping is therefore used as an algorithmic representation sensitivity, without claiming that it recovers a chemically authoritative parent structure.

A further mechanistic control is required for the regression tasks. Because a predictor that always returns the training-set mean has test MSE equal to the test-set variance plus the squared train--test mean difference, directly reducing target-mean mismatch can mechanically alter RMSE difficulty even before molecular features contribute. We therefore evaluate this mean-only predictor on every paired regression partition and distinguish that benchmark-geometry effect from the performance of ridge, random-forest, and XGBoost models.

The contribution is consequently a claim-audit framework rather than a new universal splitter or a state-of-the-art molecular predictor. We test which conclusions survive exact-size target-mean-aware selection, alternative scaffold semantics, molecular-record perturbation, and stochastic model seeds, while reporting collateral changes in benchmark composition. The broader aim is to make molecular benchmark claims auditable before they guide expensive downstream modeling or experimental decisions.
